\ifja
    \chapter{自由意志の幻想とリアリティー}
    私たちは誰もが、「自分」という不動の中心があり、自らの自由な意志によって選択し、人生を切り拓いていると信じて疑わない。毎朝ベッドから起き上がる瞬間から、重大な決断を下すその瞬間に至るまで、主導権を握っているのは常に「意識的な自己」であると思い込んでいる。

    近代社会の構造もまた、この「自由意志」という強固な前提の上に構築されている。法、道徳、教育、労働—社会のあらゆるシステムは、すべての行動を個人の自由な選択の結果と見なし、その全責任をあなた自身に問うように設計されている。常識や規範から逸脱すれば、容赦のない処罰や排除というペナルティが科され、社会はその恐怖によって私たちを規律の中に縛り付け続ける。私たちは「自己責任」という見えない鎖で己を縛り、自らの選択に怯えながら生きているのだ。

    しかし、冷徹な現代の脳科学が白日の下に晒しつつある事実は、その素朴な確信を根本から叩き壊す。

    あなたが「自らの意志で決断した」と自覚する数百ミリ秒も前に、脳の無意識の物理ハードウェアはすでに行動の引き金を引いている。意識とは司令官などではなく、生体という計算機が事後的に出力した結果を眺め、「自分が決めた」と都合よく辻褄を合わせる解釈器（インタープリター）にすぎない。私たちが絶対視してきた「自律的な意志」など、進化の過程で脳が捏造した精巧な錯覚—すなわち、美しくも残酷な幻想なのだ。
\else
    \chapter{The Illusion and Reality of Free Will}
    We all believe, without a shred of doubt, that there is an unshakeable center called the "self"—that we make choices of our own free will and carve out our own paths in life. From the moment we get out of bed each morning to the instant we make critical decisions, we are convinced that the "conscious self" is always in the driver's seat.

    The very architecture of modern society is built upon this unyielding premise of "free will." Law, morality, education, and labor—every societal system is engineered to regard all actions as the fruits of individual choice, demanding that you bear sole responsibility for them. Stray from conventional norms or common sense, and you are met with penalties: unyielding punishment and social ostracism. Society uses that terror to keep us tethered to discipline. Bound by the invisible chains of "personal responsibility," we live in perpetual fear of our own choices.

    Yet the stark realities brought to light by modern neuroscience tear this naive conviction to pieces.

    Hundreds of milliseconds before you consciously register having "made a decision of your own volition," the brain's unconscious physical hardware has already pulled the trigger on the action. Consciousness is not the commanding general; it is merely an interpreter. It gazes upon the post-hoc outputs generated by the biological computer, conveniently retrofitting a narrative to claim, "I decided this." The "autonomous will" we have held so sacrosanct is nothing more than an intricate illusion fabricated by the brain throughout evolution—a beautiful, yet merciless, fantasy.
\fi

\ifja
    \section{生体ハードウェアの実験事実：意識の遅延とモジュール構造}
    近代以前の認識論は、外界の客観的実在が感覚器を通して忠実に心に写し取られるという「素朴実在論」を信じてきた。しかし現代の神経科学が暴き出したのは、「脳は世界を受容しているのではなく、自らの仮説に基づいて世界を能動的に捏造している」という事実である。
\else
    \section{Experimental Facts of Biological Hardware: The Delay of Consciousness and Modular Architecture}
    Pre-modern epistemology placed its faith in "naive realism"—the belief that the objective reality of the external world is faithfully mirrored into the mind through our sensory organs. Yet, what modern neuroscience has laid bare is the reality that the brain does not passively receive the world; rather, it actively fabricates it based on its own hypotheses.
\fi

\ifja
    \subsection{分離脳研究と「インタープリター」モジュール}
    20世紀後半、マイケル・ガザニガ（Michael Gazzaniga）らの分離脳研究は、「一貫した自己」という近代的主体像を根底から解体した。

    難治性てんかん治療のために脳梁を切断された患者に対し、右脳のみに「歩け」という視覚指示を瞬間提示する。患者は指示通り立ち上がり、部屋を出ようとする。そこで、言語野を司る左脳に「なぜ歩き出したのですか？」と尋ねる。右脳への指示内容を知る由もない左脳は、しかし「わかりません」とは決して答えない。
    何食わぬ顔で「喉が渇いたので、コーラを買いに行くところです」と即座に辻褄の合う作話（Confabulation）を吐き出すのである。

    脳は並列分散処理を行う無数の機能モジュールの集合体であり、左脳の特定機構――ガザニガが\textbf{インタープリター（解釈器）}と呼んだモジュール――は、他モジュールが下した無意識の出力を事後的に観察し、破綻のない物語をリアルタイムで捏造しているにすぎない。
\else
    \subsection{Split-Brain Research and the `Interpreter'' Module} In the latter half of the twentieth century, split-brain research pioneered by Michael Gazzaniga and his colleagues fundamentally dismantled the modern notion of a `unified self.''

    To a patient whose corpus callosum had been severed to treat intractable epilepsy, the visual command `Walk'' was flashed exclusively to the right hemisphere. True to the instruction, the patient stood up and began to leave the room. At that moment, the left hemisphere—which houses the language faculties—was asked: `Why did you start walking?'' Having no access to the instruction presented to the right hemisphere, the left hemisphere nevertheless never replied, `I don't know.'' Without batting an eye, it instantly produced a coherent confabulation: `I'm thirsty; I'm going to get a Coke.''

    The brain is an ensemble of countless functional modules executing parallel, distributed processing. A specific mechanism in the left hemisphere—a module Gazzaniga termed the \textbf{interpreter}—merely observes the unconscious outputs yielded by other modules post hoc, fabricating an unyielding narrative in real time.    
\fi
